% the published 2021 edition of Voight, chapter 43.
\paragraph{A level bound for the non-CM fibers}
\label{sec:h5-level-repair}

This argument allows the scalar character order and coefficient field to
vary. It does not assert bounded same-field descent to a rational
GL$_2$-type variety. Let $s$ be a rational point outside the five contacts,
let $L/\mathbb Q$ be the normal fiber field constructed above, and let
$A/L$ be a non-CM fiber surface. Thus $[L:\mathbb Q]\mid576$, all
geometric homomorphisms between conjugate fibers are defined over $L$,
and $\operatorname{End}^0(A_{\overline{\mathbb Q}})=B$, the quaternion
division algebra of discriminant $210$. Put
\[
 R=B_*M(s),\qquad T=\{p:p\mid R\}.
\]
The geometric argument gives that $L$ is unramified outside $T$ and
$A$ has good reduction outside the primes above $210$. All conjugate
QM structures lie in the same normalizer orbit.

\paragraph{Good reduction after full torsion.}
First a QM surface with the indicated division action has potentially
good reduction at every finite place. Indeed, pass to an extension with
semistable reduction, using \cite[Theorem 3.4, p.~406]{silverbergzarhin}.
Every endomorphism extends uniquely to the N\'eron model by its mapping
property. On the identity component of the geometric special fiber it
preserves the maximal torus $\mathcal T$: this torus is the maximal
connected affine subgroup, and every homomorphism from a torus to an
abelian variety is zero. If its rank $t$ is positive, restriction gives
a unital homomorphism
\[
 B\longrightarrow\operatorname{End}^0(\mathcal T)=M_t(\mathbb Q),
 \qquad 1\le t\le2.
\]
The kernel is a two-sided ideal of the simple algebra $B$, and the map
is unital, so it is injective. For $t=1$ dimensions prohibit this;
for $t=2$ it would be an isomorphism from a division algebra to
$M_2(\mathbb Q)$, also impossible. Hence $t=0$. The semistable identity
component is abelian, which is good reduction. These N\'eron-model and
reduction conventions are those of \cite[\S3, pp.~405--406]{silverbergzarhin}.

Set $E=L(A[11])$. We next justify that $E/\mathbb Q$ is Galois.
For each $\sigma$, the normalizer relation supplies a compatible
isogeny ${}^{\sigma}A\to A$ of degree prime to $11$. Here is the
prime-to-$11$ detail. In the analytic description
$A_z=\mathbb C^2/(\mathcal O(z,1)^t)$ of the published 2021 edition
of Voight \cite[equations (43.1.2)--(43.1.3)]{voight2021}, a positive normalizer element
$b=\left(\begin{smallmatrix}a&b'\\c&d\end{smallmatrix}\right)$
gives the scalar compatible map $(cz+d)I:A_{bz}\to A_z$.
When $\mathcal O b\subset\mathcal O$, its degree is
$[\mathcal O:\mathcal O b]=\operatorname{Nrd}(b)^2$.
At $11$ the algebra is split and the maximal-order normalizer is
$\mathbb Q_{11}^*\operatorname{GL}_2(\mathbb Z_{11})$.
Rescale $b$ rationally to make it an $11$-adic unit, then clear the
remaining denominators by an integer prime to $11$. This produces
the claimed isogeny. Changes of level marking use isomorphisms.
All these geometric homomorphisms are $L$-defined by construction.
Consequently the $11$-torsion representations of $A$ and its conjugates
over $L$ have the same kernel. Thus every conjugate of $E$ equals $E$.

The faithful action on $A[11]$ now gives
\[
 [E:\mathbb Q]\le576|\operatorname{GL}_4(\mathbb F_{11})|
 <2^{10}11^{16}<2^{66}.
\tag{LR1}
\]
For residue characteristic different from $11$, full rational
$11$-torsion gives semistable reduction over $E$ by
\cite[Theorem 3.5, p.~406]{silverbergzarhin}. The preceding torus
argument makes this good reduction. In particular it applies at
each of $2,3,5,7$.

\paragraph{The scalar lift with controlled support.}
Choose $B$-compatible isogenies $\mu_\sigma:{}^\sigma A\to A$ over $L$,
with $\mu_1=1$, and inflate this finite choice to $G_{\mathbb Q}$.
Their discrepancies $c(\sigma,\tau)$ commute with $B$, so lie in
$\mathbb Q^*$. Put $a_\sigma=(\deg\mu_\sigma)^{-1/4}>0$ and
$\nu_\sigma=a_\sigma\mu_\sigma$, with trivial Galois action on the
added algebraic coefficients. Since a scalar $c$ has degree $c^4$,
the discrepancies of the $\nu_\sigma$ form a sign cocycle
$\varepsilon\in Z^2(G_{\mathbb Q},\{\pm1\})$, inflated from
$\operatorname{Gal}(L/\mathbb Q)$.

The sign class is unramified, hence zero, locally outside $T$.
At each odd prime where its Brauer invariant is $1/2$, choose an odd
Dirichlet character modulo $p$; at $2$ use the odd character modulo
$4$ if required. Their product $\chi$ is tame at every odd prime,
has local conductor at most $4$ at $2$, and has order
\[
 m\le2\prod_{\substack{p\in T\\p\ne2}}(p-1)\le2R.
\]
For constant coefficient groups, the boundary in
$1\to\{\pm1\}\to\mu_\infty\xrightarrow{2}\mu_\infty\to1$
is the obstruction to a character square root. Local reciprocity
identifies this obstruction with the value at $-1$: the cyclic
torsion factor has precisely this obstruction, while the valuation
and principal-unit factors admit character square roots.
Thus the boundary of $\chi$ equals $[\varepsilon]$ at every finite
place. Brauer reciprocity and the global-to-local injection give
equality globally. A pointwise square root, corrected by a sign
cochain, gives
\[
 b:G_{\mathbb Q}\longrightarrow\mu_{2m},\qquad
 b^2=\chi,\qquad\delta b=\varepsilon.
\tag{LR2}
\]
Only this \emph{sign} cochain is asserted to take root-of-unity values;
the full scalar splitting also involves the $a_\sigma$.

On $G_L$, $b$ is a conjugation-invariant character. Outside $T$ its
inertia characters are quadratic. Only finitely many are nontrivial;
twist by the quadratic character of the product of those odd primes
to kill them. At each odd prime there is a unique nontrivial quadratic
inertia character. This changes no square or coboundary in (LR2),
and introduces possible ramification only at $2\in T$.
The kernel field $L'$ of $b|G_L$ is normal over $\mathbb Q$ and
\[
 [L':\mathbb Q]\le2m[L:\mathbb Q]\le2304R.
\tag{LR3}
\]

